% !TEX root = ../nc-cubic-non-as.tex
\section{厳密計算の入力・方法・出力}
\label{sec:computation}

本付録の計算は Python 3 と SymPy を用いる有理数体上の厳密計算であり，
浮動小数点近似を用いない．

\subsection{判別曲線の滑らかさ}

次数2の入力は式\eqref{eq:q2}の determinant $\det q_2$ である．
$s_i=t_j=1$，$(i,j)\in\{0,1\}^2$ とする4チャートで，
\[
 (s_0,s_1)=
 \begin{cases}(1,x)&i=0,\\(x,1)&i=1,\end{cases}
 \qquad
 (t_0,t_1)=
 \begin{cases}(1,y)&j=0,\\(y,1)&j=1\end{cases}
\]
を同時代入して $F_{ij}\in\mathbb Q[x,y]$ を得る．
次数1の入力は式\eqref{eq:q1}の determinant $\det q_1$ であり，
$u=1,v=1,w=1$ の各チャートで残る2変数の多項式 $F$ を得る．
それぞれ
\[
 I=(F,\partial_xF,\partial_yF)
\]
の辞書式順序の Gröbner 基底を計算すると，出力はすべて $\{1\}$ である．
従って各 $I$ は単位イデアルであり，$\mathbb C$ 上にも共通零点を持たない．
7チャートが曲面を被覆するので，これは判別曲線全体の滑らかさを証明する．

\subsection{Hilbert 級数と Euler pairing}

スクリプトは偶奇の Hilbert 関数を $m$ の多項式として展開し，
本文の式との差が零多項式であることを確認する．
Hilbert 級数の比較も有理関数の差を通分・約分し，
恒等的に零になることを確認する．従って有限個の係数の一致からの推測ではない．
また式\eqref{eq:Euler2}，\eqref{eq:Euler1}の determinant と，
rank 1 の自己 Euler 数，式\eqref{eq:orbit-coordinates}の
全整数 $n$ に関する多項式恒等式を確認する．

\subsection{Hochschild homology の数値的検証}

第\ref{sec:hochschild}節の証明に現れる種数，Euler 数，
Betti 数，Hochschild homology の次元を整数として確認する．
別に Chern class の積分によって Euler 数を計算する．
次数2では
\[
 A^\bullet(\mathbb P^2\times\mathbb P^1\times\mathbb P^1)_{\mathbb Q}
 =\mathbb Q[h,a,b]/(h^3,a^2,b^2),\qquad [X_2]=2h+a+b
\]
において
\[
 c(TX_2)=\frac{(1+h)^3(1+a)^2(1+b)^2}{1+2h+a+b}\bigg|_{X_2},
 \qquad
 c_3(TX_2)=2abh-ah^2-bh^2.
\]
$[X_2]c_3(TX_2)$ の $h^2ab$ の係数は $2$ である．
次数1の cubic threefold については
\[
 c(TY)=\frac{(1+h)^5}{1+3h}\bigg|_Y,\qquad
 c_3(TY)=-2h^3,\qquad \int_Yh^3=3,
\]
従って $e_{\mathrm{top}}(Y)=-6$ となる．直線に沿う blowup で
$e_{\mathrm{top}}(\mathbb P^1)=2$ が加わり，
$e_{\mathrm{top}}(X_1)=-4$ を得る．
これは conic bundle のファイバーによる Euler 数の計算と一致する．

\Needspace{75mm}
\subsection{再現方法}

repository の root から
\begin{verbatim}
python3 -m pip install -r requirements.txt
python3 scripts/verify_all.py
python3 scripts/build.py --papers
\end{verbatim}
を実行する．検証対象は次の3ファイルである．
\begin{verbatim}
scripts/del_pezzo_degree2_verification.py
scripts/nc_dp1_quintic_checks.py
scripts/hochschild_verification.py
\end{verbatim}
この三つに加え，検証の共通入口は同梱の次数三の二つのスクリプトも実行する．
成功時の出力はその場で再生成し，重複する実行ログは収録しない．
PDF の組版には LuaLaTeX，BibTeX と日本語フォントを用いる．
必要な環境と Cursor からの実行方法は \texttt{README.md} に記す．

